\documentclass[12pt,a4paper,titlepage]{article}
\usepackage[utf8]{inputenc}
% \usepackage[francais]{babel}

\usepackage[left=1.5cm,right=2cm,top=2cm,bottom=2cm]{geometry}
\usepackage{amsmath}
\usepackage{amssymb}
\usepackage{amsthm}
\usepackage{amsfonts}
\usepackage[colorlinks]{hyperref}
\usepackage{graphicx}

% --- CSL part
\usepackage{natbib}
%\input{/project/pvs/doc/prelude}
\def\labelitemii{$\circ$}
\def\labelitemiii{$\star$}
\def\labelitemiv{$\diamond$}
\newcommand{\tcc}{{\small\small TCC}}
\newcommand{\tccs}{\tcc s}

%\renewcommand{\memo}[1]{\mbox{}\par\vspace{0.25in}\noindent\fbox{\parbox{.95\linewidth}{{\bf Memo: }#1}}\vspace{0.25in}}

\newcommand{\eg}{{\em e.g.\/},}
\newcommand{\ie}{{\em i.e.\/},}

\newcommand{\pvs}{PVS}

\newcommand{\ch}{\choice}
\newcommand{\rsv}[1]{{\rm\tt #1}}

\newcommand{\lpvstheory}[3]{\figurehead{\hozline\smaller\smaller\begin{alltt}}%
                           \figuretail{\end{alltt}\vspace{-0in}\hozline}%
                           \figurelabel{#3}\figurecap{#2}%
                           \begin{longfigure}\input{#1}\end{longfigure}}

\newcommand{\pvstheory}[3]
  {\begin{figure}[htb]\begin{boxedminipage}{\textwidth}%
   {\smaller\smaller\begin{alltt} \input{#1}\end{alltt}}\end{boxedminipage}%
   \caption{#2}\label{#3}\end{figure}}

\newcommand{\bpvstheory}[3]
  {\begin{figure}[b]\begin{boxedminipage}{\textwidth}%
   {\smaller\smaller\begin{alltt} \input{#1}\end{alltt}}\end{boxedminipage}%
   \caption{#2}\label{#3}\end{figure}}

\newcommand{\spvstheory}[1]
  {\vspace{0.1in}\par\noindent\begin{boxedminipage}{\textwidth}%
   {\smaller\smaller\begin{alltt} \input{#1}\end{alltt}}\end{boxedminipage}\vspace{0.1in}%
   }
%\newenvironment{spvstext}%
%  {\vspace{0.1in}\par\noindent\begin{boxedminipage}{\textwidth}%
%   \smaller\smaller\begin{alltt}}%
%  {\end{alltt}\end{boxedminipage}\vspace{0.1in}%
%   }


%  {\begin{boxedminipage}{\textwidth}{\smaller\smaller\begin{alltt}#1\end{alltt}}\end{boxedminipage}}
%\newenvironment{spvstheory}{\par\noindent\begin{boxedminipage}{\textwidth}\smaller\smaller\begin{alltt}}{\end{alltt}\end{boxedminipage}}

\newenvironment{pvsex}%
  {\setlength{\topsep}{0in}\smaller\begin{alltt}}%
  {\end{alltt}}

\newcommand{\pvsbnf}[2]
  {\begin{figure}[htb]\begin{boxedminipage}{\textwidth}%
   \input{#1}\end{boxedminipage}\caption{#2}\label{#1}\end{figure}}

\newcommand{\spvsbnf}[1]
  {\begin{boxedminipage}{\textwidth}\input{#1}\end{boxedminipage}}

\newcommand{\pidx}[1]{{\rm #1}} % primary index entry
\newcommand{\sidx}[1]{{\rm #1}} % secondary index entry
\newcommand{\cmdindex}[1]{\index{#1@\cmd{#1}}}
\newcommand{\icmd}[1]{\cmd{#1}\cmdindex{#1}}
\newcommand{\iecmd}[1]{\ecmd{#1}\cmdindex{#1}}
\newcommand{\buf}[1]{\texttt{#1}}
\newcommand{\ibuf}[1]{\buf{#1}\index{#1 buffer@\buf{#1} buffer}\index{buffers!\buf{#1}}}

\newenvironment{pvscmds}%
  {\par\noindent\smaller%
   \begin{tabular*}{\textwidth}{|l@{\extracolsep{\fill}}l@{\extracolsep{\fill}}l|}\hline%
     {\it Command} & {\it Aliases} & {\it Function}\\ \hline}%
  {\hline\end{tabular*}\vspace{0.1in}}

\newenvironment{pvscmdsna}%
  {\par\noindent\smaller%
   \begin{tabular*}{\textwidth}{|l@{\extracolsep{\fill}}l|}\hline%
     {\it Command} & {\it \,\,Function}\\ \hline}%
  {\hline\end{tabular*}\vspace{0.1in}}

\newcommand{\cmd}[1]{{\tt #1}}
\newcommand{\ecmd}[1]{{\tt M-x #1}}

\newcommand{\latex}{\LaTeX}                  %  LaTeX
\newcommand{\sun}{{S{\smaller\smaller UN}}}                 %  Sun
\newcommand{\sparc}{{S{\smaller\smaller PARC}}}             %  Sparc
\newcommand{\sunos}{{S{\smaller\smaller UN}OS}}             %  SunOS
\newcommand{\solaris}{{\em Solaris\/}}        %  Solaris
\newcommand{\sunview}{{S{\smaller\smaller UN}V{\smaller\smaller IEW}}} %SunView
\newcommand{\unix}{{U{\smaller\smaller NIX}}}               %  Unix
\newcommand{\lisp} {{\sc Lisp}}              %  Lisp
\newcommand{\gnu}{{Gnu Emacs}}           %  Gnu Emacs
\newcommand{\gnuemacs}{{Gnu Emacs}}      %  Gnu Emacs
\newcommand{\emacsl}{{Emacs-Lisp}}       %  Emacs Lisp
\newcommand{\shell}{{\sc Csh}}               %  C-shell

\newcommand{\update}[3]{#1\{#2\leftarrow #3\}}
\newcommand{\interp}[3]{\cal{M}\dlb {\tt #1 : #2 }\drb #3}
\newcommand{\myforall}[2]{(\forall{#1 .}\ #2)}
\newcommand{\myexists}[2]{(\exists{#1 .}\ #2)}
\newcommand{\mth}[1]{$ #1 $}
\newcommand{\labst}[2]{(\lambda{#1}.\ #2)}
\newcommand{\app}[2]{(#1\ #2)}
\newcommand{\problem}[1]{{\bf Exercise: } {\em #1}}
\newcommand{\rectype}[1]{[\# 1 \#]}
\newcommand{\recttype}[1]{{\tt [\# 1 \#]}}
\newcommand{\dlb}{\lbrack\!\lbrack}
\newcommand{\drb}{\rbrack\!\rbrack}
\newcommand{\cross}{\times}
\newcommand{\key}[1]{{\tt #1}}
\newcommand{\keyindex}[1]{\index{#1@\key{#1}}}
\newcommand{\ikey}[1]{\key{#1}\keyindex{#1}}
\newcommand{\keyword}[1]{{\smaller\texttt{#1}}}

\newenvironment{keybindings}%
  {\begin{center}\begin{tabular}{|l|l|}\hline Key & Function\\ \hline}%
  {\hline\end{tabular}\end{center}}
\def\rmif{\mbox{\bf if\ }}
\def\rmiff{\mbox{\bf \ iff \ }}
\def\rmthen{\mbox{\bf \ then }}
\def\rmelse{\mbox{\bf \ else }}
\def\rmend{\mbox{\bf end}}
\def\rmendif{\mbox{\bf \ endif}}
\def\rmotherwise{\mbox{\bf otherwise}}
\def\rmwith{\mbox{\bf \ with\ }}
\def\mapb{\char"7B\char"7B}
\def\mape{\char"7D\char"7D}
\def\setb{\char"7B}
\def\sete{\char"7D}

% ---------------------------------------------------------------------
% Macros for little PVS sessions displayed in boxes.
%
% Usage: (1) \setcounter{sessioncount}{1} resets the session counter
%
%        (2) \begin{session*}\label{thissession}
%             .
%              < lines from PVS session >
%             .
%            \end{session*}
%
%            typesets the session in a numbered box in ALLTT mode.
%
%  session instead of session* produces unnumbered boxes
%
%  Author: John Rushby
% ---------------------------------------------------------------------
\newlength{\hsbw}
\newenvironment{session}{\begin{flushleft}
 \setlength{\hsbw}{\linewidth}
 \addtolength{\hsbw}{-\arrayrulewidth}
 \addtolength{\hsbw}{-\tabcolsep}
 \begin{tabular}{@{}|c@{}|@{}}\hline 
 \begin{minipage}[b]{\hsbw}
% \begingroup\small\mbox{ }\\[-1.8\baselineskip]\begin{alltt}}{\end{alltt}\endgroup\end{minipage}\\ \hline
 \begingroup\sessionsize\vspace*{1.2ex}\begin{alltt}}{\end{alltt}\endgroup\end{minipage}\\ \hline
 \end{tabular}
 \end{flushleft}}
\newcounter{sessioncount}
\setcounter{sessioncount}{0}
\newenvironment{session*}{\begin{flushleft}
 \refstepcounter{sessioncount}
 \setlength{\hsbw}{\linewidth}
 \addtolength{\hsbw}{-\arrayrulewidth}
 \addtolength{\hsbw}{-\tabcolsep}
 \begin{tabular}{@{}|c@{}|@{}}\hline 
 \begin{minipage}[b]{\hsbw}
 \vspace*{-.5pt}
 \begin{flushright}
 \rule{0.01in}{.15in}\rule{0.3in}{0.01in}\hspace{-0.35in}
 \raisebox{0.04in}{\makebox[0.3in][c]{\footnotesize \thesessioncount}}
 \end{flushright}
 \vspace*{-.57in}
 \begingroup\small\vspace*{1.0ex}\begin{alltt}}{\end{alltt}\endgroup\end{minipage}\\ \hline 
 \end{tabular}
 \end{flushleft}}
\def\sessionsize{\footnotesize}
\def\smallsessionsize{\small}
\newenvironment{smallsession}{\begin{flushleft}
 \setlength{\hsbw}{\linewidth}
 \addtolength{\hsbw}{-\arrayrulewidth}
 \addtolength{\hsbw}{-\tabcolsep}
 \begin{tabular}{@{}|c@{}|@{}}\hline 
 \begin{minipage}[b]{\hsbw}
 \begingroup\smallsessionsize\mbox{ }\\[-1.8\baselineskip]\begin{alltt}}{\end{alltt}\endgroup\end{minipage}\\ \hline 
 \end{tabular}
 \end{flushleft}}
\newenvironment{spec}{\begin{flushleft}
 \setlength{\hsbw}{\textwidth}
 \addtolength{\hsbw}{-\arrayrulewidth}
 \addtolength{\hsbw}{-\tabcolsep}
 \begin{tabular}{@{}|c@{}|@{}}\hline 
 \begin{minipage}[b]{\hsbw}
 \begingroup\small\mbox{
}\\[-0.2\baselineskip]}{\endgroup\end{minipage}\\ \hline 
 \end{tabular}
 \end{flushleft}}
\newcommand{\memo}[1]{\mbox{}\par\vspace{0.25in}%
\setlength{\hsbw}{\linewidth}%
\addtolength{\hsbw}{-2\fboxsep}%
\addtolength{\hsbw}{-2\fboxrule}%
\noindent\fbox{\parbox{\hsbw}{{\bf Memo: }#1}}\vspace{0.25in}}
\newcommand{\nb}[1]{\mbox{}\par\vspace{0.25in}\setlength{\hsbw}{\linewidth}\addtolength{\hsbw}{-1.5ex}\noindent\fbox{\parbox{\hsbw}{{\bf Note: }#1}}\vspace{0.25in}}

%%% Local Variables: 
%%% mode: latex
%%% TeX-master: t
%%% End: 

\usepackage{boxedminipage}
\usepackage{tabularx}


\newcommand{\classnm}[1]{\texttt{\textup{#1}}}
\newcommand{\classhln}[1]{\hyperlink{#1}{\classnm{#1}}}
\newcommand{\classidx}[1]{\classnm{#1}\index{#1@\classnm{#1}}}
\newcommand{\classhtgt}[1]{\hypertarget{#1}{\classnm{#1}}\index{#1@\classnm{#1}|underline}}
\newcommand{\slotnm}[1]{{\itshape\ttfamily #1}}
\newcommand{\slotidx}[1]{\slotnm{#1}\index{#1@\slotnm{#1}}}
\newcommand{\slotidxi}[2]{\slotnm{#1}\index{#2@\slotnm{#2}}}
\newenvironment{class}[4]% name, kind, superclasses, slots
{\par\noindent\begin{boxedminipage}{\textwidth}\par\noindent\classhtgt{#1} %
$\subset$ \ifx!#3!$\emptyset$\else\classnm{#3}\fi%
\ifx!#2!\hfill[\emph{class}]\else\hfill[\emph{#2 class}]\fi%
\ifx!#4!\else\newline%
\begin{tabularx}{\linewidth}{l<{\dotfill}X}
#4
\end{tabularx}\fi\par\dotfill\par}%
{\end{boxedminipage}}
\usepackage{includes/makebnf}


\usepackage[space]{grffile} % Celui-là je sais plus à quoi il sert...

% --- Partie code
\usepackage{listings}
\lstset{language=Java} % Changer eventuellement le nom du langage
\newcommand{\codeline}[1]{\texttt{#1}}

% --- Parti math. (app.)
\newtheorem{theorem}{Theorem}
\newtheorem{lemma}[theorem]{Lemma}
\newtheorem{proposition}[theorem]{Proposition}
\newtheorem{corollary}[theorem]{Corollary}
\newtheorem{definition}[theorem]{Definition}
\newtheorem{algorithm}[theorem]{Algorithm}
\newtheorem{remark}[theorem]{Remark}

% Mathematic abbreviations

\newcommand{\N}{\mathbb{N}}
\newcommand{\Z}{\mathbb{Z}}
\newcommand{\R}{\mathbb{R}}
\newcommand{\Esp}{\mathbb{E}}
%\newcommand{\E}{\mathbb{E}} % Si tu l'utilise beaucoup
\newcommand{\Prob}{\mathbb{P}}
%\newcommand{\P}{\mathbb{P}} % Si tu l'utilise beaucoup



\title{
\vspace{-3cm}
\normalsize
\begin{tabular}{p{15cm}}
ÉCOLE POLYTECHNIQUE \\
PROMOTION X2011 \\
Férey Gaspard
\end{tabular}
\vspace{3cm}
\begin{center}
\includegraphics[height=3cm]{pictures/logo.png}
\end{center}
\vspace{1cm}
\large
\begin{center}
RAPPORT DE STAGE D'OPTION SCIENTIFIQUE\\
\vspace{1cm}
{\Huge Titre}\\
\vspace{1cm}
NON CONFIDENTIEL
\end{center}
\vspace{3cm}
\normalsize
\begin{tabular}{p{6cm} p{10cm}}
Option :                  & INFORMATIQUE \\
Champ de l'option :       & Math-Informatique \\
Directeur de l'option :   & Olivier Bournez \\
Directeur de stage :      & Olivier Bournez \\
Dates du stage :          & 7 avril - 22 aout 2014\\
Nom et adresse de l'organisme :  & SRI International \\
                          & Computer Science Laboratory (CSL) \\
                          & 333 Ravenswood Avenue \\
                          & Menlo Park, CA 94025-3493 \\
                          & United States
\end{tabular}
}

\begin{document}

\maketitle
\tableofcontents
\newpage

\section{Introduction}

\section{PVS}


\section{Translating PVS to C}

We only convert a fragment of PVS.

\section{PVS Syntax}
We describe here the syntax of PVS and the objects system used to represent them in Lisp. Some slots of the classes are voluntarily omitted. For a full description of PVS parser representation, refer to \cite{PVS:api}.


\begin{bnf}
\hproduction{expr}{Expr}{
Number \\[-0.1in] & \choice &
Name \\[-0.1in] & \choice &
Expr\ \hlink{Arguments} \\[-0.1in] & \choice &
Expr\ \hlink{Binop}\ Expr \\[-0.1in] & \choice &
\hlink{Unaryop}\ Expr \\[-0.1in] & \choice &
Expr \lit{`} \setb \hlink{Id} \choice \hlink{Number} \sete \\[-0.1in] & \choice &
\lit{(} \ites{Expr}{,} \lit{)} \\[-0.1in] & \choice &
\lit{(\#} \ites{\hlink{Assignment}}{,} \lit{\#)} \\[-0.1in] & \choice &
\hlink{IfExpr} \\[-0.1in] & \choice &
\lit{LET} \ites{\hlink{LetBinding}}{,} \lit{IN} Expr  \\[-0.1in] & \choice &
Expr \lit{WHERE} \ites{\hlink{LetBinding}}{,}  \\[-0.1in] & \choice &
Expr \lit{WITH} \lit{[} \ites{\hlink{Assignment}}{,} \lit{]} \\[-0.1in] & \choice &}
\end{bnf}



%
% The supporting productions have been moved to the file
% bnf-expr-aux.tex so that the figure gets split into two
%


%
% this (bnf-expr-aux.tex) is the latter half of bnf-expr.tex
%

\begin{bnf}

\production{Expr}{
\hlink{Number} \\[-0.1in] & \choice &
\hlink{Name} \\[-0.1in] & \choice &
Expr\ \hlink{Arguments} \\[-0.1in] & \choice &
Expr\ \hlink{Binop}\ Expr \\[-0.1in] & \choice &
\hlink{Unaryop}\ Expr \\[-0.1in] & \choice &
Expr \lit{`} \brc{ \hlink{Id} \choice \hlink{Number} } \\[-0.1in] & \choice &
\lit{(} \ites{Expr}{,} \lit{)} \\[-0.1in] & \choice &
\lit{(\#} \ites{\hlink{Assignment}}{,} \lit{\#)} \\[-0.1in] & \choice &
\hlink{IfExpr} \\[-0.1in] & \choice &
\lit{LET} \ites{\hlink{LetBinding}}{,} \lit{IN} Expr  \\[-0.1in] & \choice &
Expr \lit{WHERE} \ites{\hlink{LetBinding}}{,}  \\[-0.1in] & \choice &
Expr \lit{WITH} \lit{[} \ites{\hlink{Assignment}}{,} \lit{]} \\[-0.1in] & \choice &}

\production{Number}{\ite{\hlink{Digit}}}
\production{Id}{\hlink{Letter} \ite{\hlink{IdChar}}}
\production{IdChar}{\hlink{Letter} \choice \hlink{Digit}}
\production{Letter}{\lit{A} \choice \ldots \choice \lit{Z}}
\production{Digit}{\lit{0} \choice \ldots \choice \lit{9}}

\production{Arguments}
{\lit{(} \ites{\hlink{Expr}}{,} \lit{)}}

\production{IfExpr}
{\lit{IF} \hlink{Expr} \lit{THEN} \hlink{Expr} \\[-0.1in]
& & \rep{\brc{\lit{ELSIF} \hlink{Expr} \lit{THEN} \hlink{Expr}}} \lit{ELSE} \hlink{Expr} \lit{ENDIF}}

\production{Name}{
\lit{true}  \choice \lit{false} \choice \lit{number\_field\_pred} \choice \lit{real\_pred} \\[-0.1in]
& \choice & \lit{integer\_pred} \choice \lit{integer?} \choice \lit{rational\_pred} \\[-0.1in]
& \choice & \lit{floor} \choice \lit{ceiling} \choice \lit{rem} \choice \lit{ndiv} \choice \lit{even?} \choice \lit{odd?} \\[-0.1in]
& \choice & \lit{cons} \choice \lit{car} \choice \lit{cdr} \choice \lit{cons?} \choice \lit{null} \choice \lit{null?} \\[-0.1in]
& \choice & \lit{restrict} \choice \lit{length} \choice \lit{member} \choice \lit{nth} \choice \lit{append} \choice \lit{reverse}}

\production{Binop}{
\lit{=}  \choice \lit{\char'134 =} \choice \lit{OR} \choice \lit{\char'134 /} \choice \lit{AND} \choice \lit{\&} \choice \lit{/\char'134} \\[-0.1in]
& \choice & \lit{IMPLIES} \choice \lit{=>}  \choice \lit{WHEN} \choice \lit{IFF} \choice \lit{<=>} \\[-0.1in]
& \choice & \lit{+} \choice \lit{-} \choice \lit{*} \choice \lit{/} \choice \lit{<} \choice \lit{<=} \choice \lit{>} \choice \lit{>=}}

\production{Unaryop}{ \lit{NOT}  \choice \lit{-}}

\production{Assignment}
{ \ite{\hlink{AssignArg}} \brc{\lit{:=} \choice \lit{|->}} \hlink{Expr}}

\production{AssignArg}
{\lit{(} \ites{\hlink{Expr}}{,} \lit{)} \\[-0.1in]
& \choice & \lit{`} \hlink{Id} \\[-0.1in]
& \choice & \lit{`} \hlink{Number}}

\production{LetBinding}
{\brc{\hlink{LetBind} \choice \lit{(} \ites{\hlink{LetBind}}{,} \lit{)}} \lit{=} \hlink{Expr}}

\production{LetBind}
{ \hlink{Id} \opt{\lit{:} \hlink{TypeExpr}}}

\end{bnf}


%\production{AssignArgs}
%{\hlink{Id} \opt{\lit{!} \hlink{Number}}\\[-0.1in]
%& \choice & \hlink{Number}\\[-0.1in]
%& \choice & \ite{\hlink{AssignArg}}}

%\production{IdOp}
%{\hlink{Id} \ch \hlink{Opsym} \ch \hlink{Number}}
%
%\production{Opsym}
%{\hlink{Binop} \choice \hlink{Unaryop} \choice \lit{TRUE}
%\choice \lit{FALSE}}


%\production{BindingOp}
%{\lit{LAMBDA} \choice \lit{FORALL} \choice \lit{EXISTS} \choice \{ \hlink{IdOp} \lit{!} \}}

%\production{IdOps}
%{\ites{\hlink{IdOp}}{,}}

%\production{Selection}
%{\hlink{IdOp} \opt{\lit{(} \hlink{IdOps} \lit{)}} \lit{:} \hlink{Expr}}




\section{Types}
A PVS theory can be typechecked using the emacs interface \codeline{M-x typecheck} or with Lisp function \codeline{(tc name-theory)}. This first runs the PVS parser on the code and generates CLOS objects to represent it. Then, the PVS typechecker is run on this internal representation of the theory and tries to give a type to all expressions generating TCC when needed.\\

Here we describe how PVS types are represented in Lisp. The syntax of PVS we allow


\begin{bnf}

\hproduction{type-expr}{TypeExpr}
{\hlink{Name} \\[-0.1in]
& \choice & \hlink{EnumerationType} \\[-0.1in]
& \choice & \hlink{Subtype} \\[-0.1in]
& \choice & \hlink{TypeApplication} \\[-0.1in]
& \choice & \hlink{FunctionType} \\[-0.1in]
& \choice & \hlink{TupleType} \\[-0.1in]
& \choice & \hlink{CotupleType} \\[-0.1in]
& \choice & \hlink{RecordType}}

\hproduction{enumtype}{EnumerationType}
{\lit{\setb} \hlink{IdOps} \lit{\sete}}

\hproduction{subtype}{Subtype}
{\lit{\setb} \hlink{SetBindings} \lit{|} \hlink{Expr} \lit{\sete} \\[-0.1in]
& \choice & \lit{(} \hlink{Expr} \lit{)} }

\hproduction{type-application}{TypeApplication}
{\hlink{Name}\ \hlink{Arguments}}

\hproduction{funtype}{FunctionType}
{\opt{\lit{FUNCTION} \choice \lit{ARRAY}} \\[-0.1in]
& & \lit{[} \ites{\setb \opt{\hlink{IdOp} \lit{:}} \hlink{TypeExpr} \sete}{,} \lit{->} \hlink{TypeExpr} \lit{]}}

\hproduction{tupletype}{TupleType}
{\lit{[} \ites{\setb \opt{\hlink{IdOp} \lit{:}} \hlink{TypeExpr} \sete}{,} \lit{]}}

\hproduction{cotupletype}{CotupleType}
{\lit{[} \ites{\setb \opt{\hlink{IdOp} \lit{:}} \hlink{TypeExpr} \sete}{+} \lit{]}}

\hproduction{recordtype}{RecordType}
{\lit{[\#} \ites{\hlink{FieldDecls}}{,} \lit{\#]}}

\hproduction{field-decl}{FieldDecls}
{\hlink{Ids}\ \lit{:} \hlink{TypeExpr}}
\end{bnf}


\begin{class}{type-expr}{abstract}{\classhln{syntax}}
{}
\end{class}
\begin{class}{type-name}{}{\classhln{type-expr} \classhln{name}}
{\slotidx{adt} & }
\end{class}
\begin{class}{subtype}{}{\classhln{type-expr}}
{\slotidx{supertype} & \\
 \slotidx{predicate} & }
\end{class}
\begin{class}{funtype}{}{\classhln{type-expr}}
{\slotidx{domain} & \\
 \slotidx{range} & }
\end{class}
\begin{class}{tupletype}{}{\classhln{type-expr}}
{\slotidx{types} & }
\end{class}
\begin{class}{recordtype}{}{\classhln{type-expr}}
{\slotidx{fields} & }
\end{class}




boolean, number, number\_field, real, 
rational, integer, $A \rightarrow B$, restricted types
$\text{below}(10) := \{ x:\text{int} | 0 \leq x < 10 \} $)  \\
enum
datatype

C types:
[unsigned] char, int, long, double
boolean
arrays
strings
enum
struct
and others: short int, float, union, size\_t, ...



\subsection{Fragment of PVS syntax}





\subsection{Difficulties}
if-expr
update-expr





\section{Other activities at SRI}

Robin project, HACMS \\
Contest week-end 14-15 June \\
Summer School \\






\section{Other works at SRI}
Discovering PVS :
Translating Coq proofs to PVS
PVS library for basic linear algebra


HACMS with Robin
Parsing Lisp code -> generate HTML architecture file
Correcting translator PVS to SMT-LIB

Summer school






Try using bibtex here

-- PVS --

PVS API Reference
PVS Lisp sources (github rep)
PVS language Reference
PVS System Guide
PVS Prelude library

-- Lisp --
Common Lisp Guy L. Steele Jr

-- C-- 
The C Library Reference Guide, Eric Huss

-- Other --
Compilation course, J.C. Filliatre


\cite{PVS:language}
\cite{PVS:manuals}
\cite{PVS:userguide}


\bibliographystyle{plain}
\bibliography{report}


\end{document}
